
@@M0@@ est un @@M1@@-idéal.

Soit alors @@M2@@ un élément d'augmentation @@M3@@ de @@M4@@~; @@M5@@ est congru @@M6@@ à @@M7@@, et par suite, @@M8@@ est congru à @@M9@@, d'où :
@@M10@@

\textbf{Corollaire 6.6.4.} \uline{Soit} @@M11@@ \uline{l'homomorphisme d'augmentation de} @@M12@@~; \uline{alors}, \uline{pour tout} @@M13@@, \uline{on a} :
@@M14@@

Par définition, @@M15@@. Comme @@M16@@ est d'augmentation @@M17@@ on peut lui appliquer le corollaire 6.6.2, d'où le résultat.

\textbf{Corollaire 6.6.5.} \uline{Avec les notations de} 3.10, @@M18@@.

D'après le corollaire précédent, @@M19@@. Le corollaire résulte alors de ce que les deux filtrations sont des filtrations d'anneau, et de ce que les @@M20@@, pour @@M21@@ et @@M22@@ variables, engendrent la filtration de @@M23@@-anneau.

\textbf{Corollaire 6.6.7.} \uline{Soit} @@M24@@ \uline{une} @@M25@@-\uline{algèbre graduée de caractéristique nulle}, \uline{à degrés bornés}~; \uline{alors les deux filtrations définies sur} @@M26@@ \uline{sont} \uline{identiques}.

Supposons que @@M27@@, et soit @@M28@@ un élément de @@M29@@. Posons @@M30@@~; on a alors :
@@M31@@, car @@M32@@ est d'augmentation nulle. Par suite, 
